\noindent\textbf{Por qué se usa este procedimiento.} Integrar los flujos de las alas da un par respecto al alma. Igualar ese par a $C_y e_c$ y usar $I_z=th^2(h+6b)/12$ elimina $t$ y permite despejar $e_c=3b^2/(h+6b)$. La distancia se mide desde el alma.

\begin{enumerate}

\item No depende del espesor.

\end{enumerate}

\noindent\textbf{Operaciones y resultados intermedios.}

\begin{enumerate}

\item Distancia al alma.

\[\fitmath{y_{1}=\frac{3\cdot {60}^{2}}{120+6\cdot 60}\approx 22.5\,\mathrm{mm}}\]

\end{enumerate}
